\scanpage{53}

\mspage{- 133 -}\hfill n$^\circ$371

$n-2$ \uline{de} $\underline X$ \uline{obtenu en supprimant dans} $\underline X$ \uline{les lignes d'indices} $i$ \uline{et} $j$ \uline{et les colonnes
d'indices} $h$ \uline{et} $k$ ; \uline{on a alors}
\[
(38)\qquad X^{ih}X^{jk}-X^{ik}X^{jh}=(\det\underline X).X^{ij,hk}.
\]
Il suffit d'appliquer la formule (35) lorsque $H=I-\{i,j\}$, $K=I-\{h,k\}$,
où $I=\unsure{[\![1,n]\!]|[1,n]}$.

Nous verrons au chap.IV que les formules (34) à (38) sont encore
valables sans aucune hypothèse sur la matrice carrée $\underline X$.

\begin{center}\texttt{-:-:-:-:-}\end{center}

\margin{\rotatebox{180}{52}}
